\documentclass[9pt,twocolumn,twoside]{pnas-new}
% Use the lineno option to display guide line numbers if required.

\usepackage[T1]{fontenc} % Icelandic letters (þ, ð)
\usepackage{xcolor}
\newcommand{\TODO}[1]{\textcolor{red}{[TODO: #1]}}
% Fallbacks in case this version of pnas-new.cls lacks a research-article command
\providecommand{\significancestatement}[1]{}
\providecommand{\matmethods}[1]{\section*{Materials and Methods}#1}
\providecommand{\showmatmethods}{}

\templatetype{pnasresearcharticle} % research article: has Significance Statement + Materials and Methods

\begin{document}

\title{Do feuds follow family lines? Killing and kinship in the Sagas of Icelanders}

\author[a]{Ylfa Margrét Ólafsdóttir}
\author[a]{Tanja Kristín Árnadóttir}
\author[a]{Andrea Hermóðsdóttir}

\affil[a]{Technical University of Denmark, 02805 Social Graphs and Interactions, Kongens Lyngby, Denmark}

\leadauthor{Ólafsdóttir}

\significancestatement{The Sagas of Icelanders are often read as stories of family feuds. We test this reading on all 42 sagas by building a network of every named killing, 482 in all, together with the family ties of the people involved and the motive of each killing. Most saga killings are between families and part of a feud, and a quarter are revenge, usually aimed at the killer himself. Blood relatives avenge each other and almost never kill each other, but ties made by fosterage or marriage offer no such protection. We also show how such networks can be built from historical texts with AI assistance while reporting their accuracy openly.}

\authorcontributions{\TODO{Confirm with the group. Draft based on the project plan: Y.M.Ó. built the family and killing networks and did the network analysis; T.K.Á. handled data and validation; A.H. did the bridging and language analysis. All authors defined the annotation criteria, checked samples by hand and wrote the paper.}}
\authordeclaration{The authors declare no competing interests.}
\correspondingauthor{\textsuperscript{1}To whom correspondence should be addressed. E-mail: \TODO{e-mail}}

\keywords{Icelandic sagas $|$ feud $|$ kinship networks $|$ structural balance $|$ null models}

\begin{abstract}
The Sagas of Icelanders are famous for their feuds, but whether saga violence follows family lines has not been tested across the whole corpus. We extracted every named killing from all 42 sagas, 482 killings among 624 characters, together with the family ties between killer and victim and the motive of each killing. Candidate sentences were selected by keyword, annotated by an AI model under rules we defined, recoded in a blind second pass and checked by hand. The signed network of kinship (friendly) and killing (hostile) ties is small-world and structurally balanced: 11\% of closed triads contain an odd number of hostile edges, against 42\% by chance. Compared with a null model that reshuffles victims within each saga, killings within a household or kin group occur at a third of the chance rate (3.2\% vs.\ 9.0\%), and killings between immediate blood relatives almost never happen (2 vs.\ 14 expected). Ties made by fosterage or marriage, by contrast, offer no protection: they join killer and victim exactly as often as chance predicts. A quarter of all killings are revenge; avengers are mostly blood kin of the dead, and revenge strikes the original killer himself in two thirds of cases. In the sagas, blood binds: blood kin avenge each other and almost never kill each other, while the made ties of fosterage and marriage are crossed by violence as readily as no tie at all.
\end{abstract}

\dates{This manuscript was compiled on \today}
\doi{\url{www.pnas.org/cgi/doi/10.1073/pnas.XXXXXXXXXX}}

\maketitle
\thispagestyle{firststyle}
\ifthenelse{\boolean{shortarticle}}{\ifthenelse{\boolean{singlecolumn}}{\abscontentformatted}{\abscontent}}{}

\dropcap{T}he Sagas of Icelanders (\textit{Íslendingasögur}) are 42 prose narratives, written in the 13th and 14th centuries, about Icelanders of the 9th to 11th centuries. Their best-known theme is the feud: killings answered by killings, often between families \cite{byock1982,miller1990}. This project investigates relationships of lethal violence in the sagas. We construct a directed network in which nodes are named human characters and an edge points from a killer to their victim, and we combine it with the family relations of each participant, to ask whether killings follow family lines or cut across them.

Network analysis of the sagas has so far treated them as social networks. Mac Carron and Kenna built character networks by hand for 18 sagas, distinguished friendly from hostile links, and found that saga society looks like a real social network: small-world and structurally balanced \cite{maccarron2013}. They grouped family with other friendly ties and did not ask whether violence follows family structure. We build kinship and killing as separate layers, from the Icelandic text, for all 42 sagas; we separate kinship into blood, marriage and fosterage; and we test against null models whether killings follow family lines (SQ1), whether family communities match the groups that fight each other and which kinds of tie violence follows (SQ2), and how revenge runs through the network.

\section*{Results}

\subsection*{The dataset}
Across the 42 sagas we identified 482 distinct killings among 624 named characters: 253 killers and 449 victims, 78 of whom are both (Table~\ref{tab:overview}). All sagas but one contain at least one killing, and \textit{Njáls saga} alone has 90. Most killers kill once (177 of 253); 18 kill five times or more, led by Kári (15), Gunnar of Hlíðarendi, Grettir and Þorgeir Hávarsson (12 each). The killing network falls into 150 chains of linked killings, the largest of which (65 people) is the conflict in \textit{Njáls saga}. Under our strict definition, 457 killings (95\%) are between people with no family tie; 15 are between blood kin, 5 between foster kin, 4 between in-laws and 1 between sworn brothers, and 36 fall within a family group or household. Counting also indirect ties (e.g.\ a sister's husband's brother), 39 killings (8\%) are between kin.

\begin{table}[tbhp]
\centering
\caption{The dataset.}
\label{tab:overview}
\begin{tabular}{lr}
\toprule
Sagas / words & 42 / 897k \\
Candidate sentences (original + extension) & 1,048 + 1,745 \\
Distinct killings & 482 \\
People (killers / victims / both) & 624 (253 / 449 / 78) \\
Killing chains (components); largest & 150; 65 people \\
Family network: people / ties & 1516 / 1650 \\
Feuds (2+ killings) / killings in feuds & 93 / 391 \\
Revenge killings (explicit in text) & 127 (79) \\
Killer and victim related (strict / wide) & 25 / 39 \\
\bottomrule
\end{tabular}
\end{table}

\subsection*{Network structure}
Combining the two layers gives a signed network in which family ties are friendly edges and killings hostile edges, the distinction used by Mac Carron and Kenna \cite{maccarron2013}. For all sagas it has 1{,}673 people and 2{,}117 edges, 479 of them hostile (Table~\ref{tab:network}). The saga networks are small-world: path lengths are close to those of random graphs of the same size and density while clustering is an order of magnitude higher (\textit{Njáls saga}: $\ell=4.3$ vs.\ $4.8$, $C=0.15$ vs.\ $0.013$, $S=14$), and they have strong community structure (Louvain $Q=0.64$--$0.74$ per saga). Unlike the full character networks of Mac Carron and Kenna, ours are disassortative ($r_k<0$), because they are built around the participants in killings and the relatives recorded for them.

The signed network is structurally balanced. Only 12\% of closed triads in \textit{Njáls saga} and 11\% in all sagas contain an odd number of hostile edges, against 41--42\% when the signs are shuffled over the same edges ($p=0.002$ for both). This is close to the 10\% that Mac Carron and Kenna report for \textit{Njáls saga} from an independently built network \cite{maccarron2013}. In saga terms, a killing rarely closes a triangle with family ties: people seldom kill someone tied to their own relatives. Fig.~\ref{fig:njala} shows the network of \textit{Njáls saga}: killings link family clusters, such as Gunnar's kin at Hlíðarendi and Njáll's household at Bergþórshvoll, rather than running inside them.

\begin{table*}[t!]
\centering
\caption{The signed saga networks: family ties as friendly edges, killings as hostile edges. $N$, $M$: nodes and edges (hostile in brackets); $\langle k\rangle$, $k_{\max}$: mean and maximum degree; $\ell$, $C$: mean path length (giant component) and clustering, with random-graph equivalents; $S$: small-world-ness; $G_C$: giant component; $r_k$: degree assortativity; $Q$: Louvain modularity; $\Delta$: closed triads with an odd number of hostile edges, with the mean of 1{,}000 sign shuffles. $\Delta$ rests on 177 triads in \textit{Njáls saga}, 518 in all sagas, and 11--39 in each smaller saga. Characters are not merged across sagas, so the all-sagas network is a union of saga networks.}
\label{tab:network}
\begin{tabular*}{\textwidth}{@{\extracolsep{\fill}}lrrrrrrrrrrrrr@{}}
\toprule
Network & $N$ & $M$ (hostile) & $\langle k\rangle$ & $k_{\max}$ & $\ell$ & $\ell_{\mathrm{rand}}$ & $C$ & $C_{\mathrm{rand}}$ & $S$ & $G_C$ & $r_k$ & $Q$ & $\Delta$ (null) \\
\midrule
\textit{Njáls saga} & 253 & 404 (90) & 3.19 & 32 & 4.26 & 4.77 & 0.154 & 0.0126 & 13.7 & 92\% & -0.12 & 0.68 & 12.4\% (41.3\%) \\
\textit{Egils saga} & 76 & 104 (21) & 2.74 & 21 & 3.47 & 4.3 & 0.246 & 0.036 & 8.5 & 93\% & -0.44 & 0.64 & 16.1\% (40.1\%) \\
\textit{Grettis saga} & 79 & 104 (23) & 2.63 & 21 & 3.51 & 4.51 & 0.203 & 0.0333 & 7.8 & 59\% & -0.21 & 0.68 & 0.0\% (42.1\%) \\
\textit{Laxdæla saga} & 83 & 103 (17) & 2.48 & 19 & 3.65 & 4.86 & 0.098 & 0.0299 & 4.4 & 86\% & -0.44 & 0.66 & 6.2\% (35.8\%) \\
\textit{Eyrbyggja saga} & 74 & 85 (14) & 2.3 & 11 & 3.75 & 5.17 & 0.117 & 0.031 & 5.2 & 88\% & -0.39 & 0.69 & 9.1\% (36.0\%) \\
\textit{Reykdæla saga} & 68 & 82 (23) & 2.41 & 13 & 3.87 & 4.79 & 0.206 & 0.0355 & 7.2 & 68\% & -0.35 & 0.74 & 0.0\% (47.0\%) \\
\textit{All 42 sagas} & 1673 & 2117 (479) & 2.53 & 32 & 4.26 & 7.99 & 0.154 & 0.0015 & 191.2 & 14\% & -0.19 & 0.95 & 10.8\% (41.9\%) \\
\bottomrule
\end{tabular*}
\end{table*}

\begin{figure*}[t!]
\centering
\includegraphics[width=17.8cm]{network_njala.pdf}
\caption{The signed network of \textit{Njáls saga} (giant component). Gray: family ties (friendly edges); red: killings (hostile edges); blue: killers and victims, sized by degree; small gray dots: relatives recorded for them. Killings run between family clusters rather than within them.}
\label{fig:njala}
\end{figure*}

\subsection*{SQ1: Killings and family lines}
Killer and victim belonged to the same family group or household in 15 of the 468 killings where both groups are known (3.2\%). Under the null model, which reshuffles victims within each saga, the expected share is 9.0\% $\pm$ 1.0\% ($z=-6.0$, $p=0.002$). The result holds without the 71 partly supported killings and when only killings where both families are known are kept (both $p=0.002$). A direct kin tie between killer and victim in the family network is also rarer than chance: 12 of 482 killings (2.5\%) against 4.9\% $\pm$ 0.8\% ($z=-3.0$, $p=0.008$). Saga characters kill outside their own family and household far more often than chance predicts.

\subsection*{SQ2: Family communities, feuds and types of tie}
The family network of the killing participants and their kin has 1{,}516 people and 1{,}650 ties. Louvain divides it into 232 communities ($Q=0.97$), stable over 100 seeds (mean pairwise NMI $1.00$) and in agreement with label propagation (NMI $0.93$); the high modularity partly reflects that the network consists of many small kin clusters. Feuds, the connected components of the killing network, line up with family communities slightly more than shuffled feuds do (NMI $0.839$ vs.\ $0.825 \pm 0.004$, $z=3.7$, $p=0.002$; blood ties only: $0.838$ vs.\ $0.828$, $p=0.014$): feuds involve a recurring set of families, but the effect is small.

The type of tie matters (Fig.~\ref{fig:ties}). Only 2 killings in the corpus are between immediate blood relatives (parent, child or sibling), where chance predicts 14.2 ($p=0.002$). Killings between in-laws (5 vs.\ 4.4, $p=0.87$) and between foster kin (5 vs.\ 4.8, $p=1.0$) occur exactly at the chance rate. The protection against being killed by one's own kin is a protection of blood; the made ties of marriage and fosterage give none.

\begin{figure}[tbhp]
\centering
\includegraphics[width=\linewidth]{sq2_ties.pdf}
\caption{SQ2. Killings between killer and victim joined by a direct tie of each type (blue), against the 95\% range (gray bar) and mean (tick) of 1{,}000 within-saga permutations. Immediate blood kin almost never kill each other; in-laws and foster kin do so at the chance rate.}
\label{fig:ties}
\end{figure}

\subsection*{Feuds and revenge}
Read for motive and context, the 482 killings form 179 conflicts. 93 are feuds with two or more linked killings, and they contain 391 killings (81\%); the largest are the killing of Höskuldur Hvítanesgoði and the Burning of Njáll (27 killings), Gunnar of Hlíðarendi's conflict with Otkell, Starkaður and their allies (19) and Gull-Þórir's feud with Steinólfur and Hof-Hallur (13). Revenge is the most common motive (125 killings), followed by self-defence (108), honour or insult (60), property disputes (50) and battle (45). In all, 127 killings (26\%) avenge an earlier death, 79 of them with an explicit statement in the text.

Revenge is aimed precisely (Fig.~\ref{fig:revenge}B). In 80 of the 127 revenge killings (63\%) the avenger kills the original killer himself; the killer's kin are struck in 25 and his household or allies in 21. A network measure computed independently from the family ties agrees: avengers kill the original killer or someone directly tied to him almost twice as often as chance (70 vs.\ 37.9 under victim reshuffling, $p=0.002$), and the text confirms 52 of those 70 cases. Avengers are mostly blood kin of the person they avenge (67 of 127), followed by in-laws (19), household members (14) and foster kin (8). Among avengers related to the avenged, this split (71\% blood, 20\% marriage, 9\% foster) is close to the split of all recorded kin ties (75\%, 21\%, 4\%; Fig.~\ref{fig:revenge}A): the duty of revenge falls on whichever kin there are.

\begin{figure*}[t!]
\centering
\includegraphics[width=17.8cm]{revenge.pdf}
\caption{Revenge, from the text-based annotation of motives (127 revenge killings). (A) The avenger's tie to the person avenged (blue, with counts), compared with the share of each tie type among all recorded kin ties (gray). (B) Whom revenge strikes.}
\label{fig:revenge}
\end{figure*}

\subsection*{Robustness}
On the original keyword set alone (194 killings found with forms of \textit{drepa}, \textit{vega} and \textit{víg}), killings within a family group are also rarer than chance (2.2\% vs.\ 5.8\%, $p=0.004$), immediate blood kin never kill each other (0 vs.\ 4.7, $p=0.008$), and the network is structurally balanced (15\% vs.\ 35\%). Two results did not survive the extension: in the smaller set, direct kin ties were as common as chance, and foster kin appeared to kill each other more often than chance (5 vs.\ 2.4, $p=0.08$). With 2.5 times as many killings, both disappear.

\section*{Discussion}
Saga violence is overwhelmingly violence between families, and most of it is organised into feuds: 95\% of killings are between people with no family tie, 81\% belong to a feud, and killings within a household or kin group occur at a third of the chance rate. In network terms, the signed network of kinship and killing is structurally balanced, as Mac Carron and Kenna found for saga character networks \cite{maccarron2013}: hostile edges run between family clusters, not through them. A quarter of all killings are revenge, aimed in two thirds of cases at the killer himself, consistent with historians' view of saga feud as regulated rather than indiscriminate violence \cite{byock1982,miller1990}.

The most revealing result concerns the kind of tie. Blood kin almost never kill each other, and they carry out most revenge. Ties made by marriage and fosterage, by contrast, give no protection at all: in-laws and foster kin kill each other as often as unrelated pairs placed at random would. Some of the most famous killings in the corpus run along exactly these ties: the sons of Njáll killing their foster-brother Höskuldur Hvítanesgoði, and Bolli killing his cousin and foster-brother Kjartan in \textit{Laxdæla saga}. Marriage and fosterage were the ties by which saga households built alliances; our results suggest that such alliances did not carry the solidarity of blood. \TODO{Group: discuss, add one or two sentences in your own words.}

\textit{Limitations.} Candidates were found by keyword, so killings told without any of our killing words are missed; the agents that verified the extension noted dozens of such killings, and a recall check on three sagas read in full found \TODO{X of Y} named killings in the dataset. Only killings with named participants are included, and an event told in two sagas is counted once, in the saga that narrates it. The family network contains only the kin recorded for the participants, and only direct ties. Family-group labels assign followers to their master's household, so part of the SQ1 result reflects that people on the same side rarely kill each other. Annotation, kinship lookup and motive coding were done by AI models; manual checks found the original killings largely correct but the kin codings weaker, and the counts of kin killings remain small. Characters are not merged across sagas.

\matmethods{
\subsection*{Corpus and candidates}
We used all 42 sagas in modern-spelling Icelandic from the Icelandic Saga Database (sagadb.org), about 0.9 million words. The original candidate set (1{,}048 sentences) contained one of 14 forms of \textit{drepa}, \textit{vega} and \textit{víg}. Because a recall check showed that many killings are told with other words, we extended it with 1{,}745 further sentences containing other killing words (e.g.\ \textit{féll}, \textit{hjó}, \textit{banahögg}, \textit{bana}, \textit{brenndur}, \textit{drepa}, \textit{drepur}). Matches are only candidates: \textit{drap á dyr} means knocking on a door, \textit{hjó} can be a blow that does not kill, and \textit{féll} has many other senses.

\subsection*{Annotation}
A positive label required an actual killing of a human by a human, killer and victim both identifiable by name from the sentence or its context, and an attribution the text supports. We excluded threats, predictions, negated killings, dreams, non-human participants (including victims killed in animal shape), death by sorcery, suicide, anonymous participants and unconfirmed attributions; a person who ordered a killing counts as killer only if present. The detailed annotation was done by AI models (Claude, Anthropic) under these rules, which we refined after inspecting their output. Repeated mentions of the same killing are merged into one directed killer$\rightarrow$victim edge; parallel lists are paired in order; and only where a saga credits a killing to a whole party does each member get an edge. Characters are identified by name within a saga. An event told in two sagas is counted once, in the saga that narrates it (16 such duplicates were removed).

\subsection*{Family situation and relations}
For every killer and victim, an AI model looked up parents, siblings, spouses, children, in-laws, foster and sworn kin in the saga text, mainly from the character's introduction (\textit{X hét maður, sonur Y}), citing a chapter for every entry and marking families the saga does not give as not found. \textit{Coding rules.} \textit{Blood}: descent or siblinghood up to second cousins. \textit{Affinity}: one marriage link existing at the time of the killing. \textit{Foster}: a direct fosterage tie. \textit{Sworn}: explicit sworn brotherhood. Longer chains, later marriages and betrothals count as no tie; a wide definition used for robustness also counts them. Killer and victim are in the same family group if they shared a household (including master and slave), were blood kin up to first cousins, or were tied by direct fosterage.

\subsection*{Verification}
For the original 194 killings, a blind second AI pass recoded every killing from the text alone; it contradicted none of the attributions, and the two passes agreed on kin vs.\ not kin in 93.9\% of killings (Cohen's $\kappa=0.74$); a third pass settled the remaining disagreements under the written rules. Each author then checked a different random sample of 30 of these killings by hand (90 in total, stratified to include all 22 coded as kin); the judgements were the authors' own, with an AI model occasionally used to interpret the older Icelandic. They confirmed 82 of 90 attributions (5 partly correct, 3 wrong) and 79 of 90 relation codings, but only 14 of 22 kin codings; 0 of 10 sampled negative labels were killings. For the 291 killings added by the extension, a separate AI pass verified each one against the text, removed duplicates and name variants, and coded family, relation and motive; \TODO{N of 30} sampled new killings were confirmed by hand. All corrections (40 in total) are documented with their textual evidence. Two AI passes can share biases, so agreement between them measures consistency, not accuracy.

\subsection*{Feuds and motives}
For each killing an AI model recorded the feud it belongs to (a chain of hostile acts the saga links causally), the motive, and, only where the text states it or makes it unmistakable, which earlier death it avenges, with an Icelandic quote, the avenger's relation to the avenged person and whom the revenge strikes. This annotation was blind to the network-based revenge measure. The first author checked 20 revenge links and 10 other killings against the text: 25 of 30 motives and all 20 revenge links were confirmed.

\subsection*{Signed network}
We combined the family graph (friendly edges) with the killings (hostile edges; a killing overrides a family tie between the same pair) and computed the measures of Mac Carron and Kenna \cite{maccarron2013}: mean path length on the giant component, clustering, their equivalents for a random graph of the same size and mean degree ($\ell_{\mathrm{rand}}=\ln N/\ln\langle k\rangle$, $C_{\mathrm{rand}}=\langle k\rangle/N$), small-world-ness $S=(C/C_{\mathrm{rand}})/(\ell/\ell_{\mathrm{rand}})$, degree assortativity, Louvain modularity, and structural balance $\Delta$, the share of closed triads with an odd number of hostile edges, compared with 1{,}000 shuffles of the signs over the same edges.

\subsection*{Null models}
SQ1 needs a family group for every person, so it uses the per-person family-group label from the family lookup (available for 97\% of participants); a second statistic uses direct ties in the family network (parent, child, sibling, spouse, in-law, foster). We permute victims among killings within each saga 1{,}000 times, which preserves each character's number of killings committed and suffered, rejecting permutations with self-killings, and report $z$-scores and empirical two-sided $p$-values. For SQ2 we ran Louvain \cite{blondel2008} with 100 seeds and compared it with label propagation \cite{raghavan2007} by normalised mutual information (NMI). Conflict groups are the connected components of the undirected killing network; because both partitions are nested within sagas, which inflates NMI, the observed NMI is compared with that of feuds rebuilt from the permutations. Killings along direct blood, affinity and foster ties are counted against the same null model. For the network revenge measure, killings were ordered by chapter; a killing counts as targeted revenge if the killer has a direct tie to an earlier victim and the victim is that victim's killer or directly tied to him. It is compared with victim reshuffling and with reshuffling the order of killings.
}

\showmatmethods{}

\dataavail{The saga texts are available from sagadb.org. Labelled sentences, all corrections, the family tables, the feud annotation, the final killing list (\texttt{killings\_all.csv}) and all scripts are available at \url{https://github.com/s252869/sagas}.}

\acknow{\TODO{e.g. We thank the course staff of 02805 for feedback.} Annotation and family lookup were done with Claude (Anthropic), as described in Materials and Methods.}

\showacknow{}

\begin{thebibliography}{9}
\bibitem{byock1982} J.~L. Byock, \textit{Feud in the Icelandic Saga} (University of California Press, Berkeley, 1982).
\bibitem{miller1990} W.~I. Miller, \textit{Bloodtaking and Peacemaking: Feud, Law, and Society in Saga Iceland} (University of Chicago Press, Chicago, 1990).
\bibitem{maccarron2013} P.~Mac~Carron, R.~Kenna, Network analysis of the Íslendinga sögur -- the Sagas of Icelanders. \textit{Eur. Phys. J. B} \textbf{86}, 407 (2013).
\bibitem{blondel2008} V.~D. Blondel, J.-L. Guillaume, R.~Lambiotte, E.~Lefebvre, Fast unfolding of communities in large networks. \textit{J. Stat. Mech.} P10008 (2008).
\bibitem{raghavan2007} U.~N. Raghavan, R.~Albert, S.~Kumara, Near linear time algorithm to detect community structures in large-scale networks. \textit{Phys. Rev. E} \textbf{76}, 036106 (2007).
\end{thebibliography}

\end{document}
